\section{制度不是脚注：它决定了仓位设计}

把全球私募方法搬到中国个人账户时，最容易遗漏的是市场微观结构。机构可以依靠协议条款、长期资本、交叉资产流动性或专门交易团队管理退出；普通公开市场投资者需要面对交易制度、价格限制、停复牌、订单簿深度、账户权限和税费。于是，稳健的研究框架首先是一种\textbf{不会假设不存在的流动性}的设计。

\begin{exhibitbox}[Exhibit 5：中国个人版方法的制度适配]
\small
\begin{tabularx}{\exhibitboxwidth}{L{3.0cm} X L{4.1cm}}
\toprule
\textbf{约束} & \textbf{框架影响} & \textbf{研究动作} \\
\midrule
普通现货的结算 / 回转交易边界 & 普通股票应按交收前不得卖出的规则设计；有限例外必须逐只核验，不能把卖出设计成一定可当日执行的止损。 & 按非日内可处置的压力情景建仓；核验具体品种、现行登记结算和交易规则。 \\
涨跌幅、停牌和上市初期不连续性 & 价格限制不是全市场恒定数字：板块、基金类别、特殊情形和新上市阶段可不同；无涨跌幅也不等于一定有流动性。 & 排除或特别降级新上市、特殊处理、长期停牌、连续涨跌停或治理事件高发对象；限价与分期不等于保证成交。 \\
融资融券的准入和维持要求 & 融资并不等于可复制杠杆收购（LBO）的稳定资本结构，融券也不等于稳定可得的对冲。 & 默认零融资加仓、零裸空头；不让借款成为策略成立条件。 \\
交易型开放式指数基金（ETF）机制差异 & ETF 的二级市场、申赎和回转安排取决于类别、产品文件和当期规则；仅有某类 ETF 的机制不能外推到权益、跨境、商品或 REIT。 & 不把“可申赎”错当成二级市场必然的即时退出；查看基金合同、交易规则和流动性。 \\
公募不动产投资信托基金（REIT）产品结构 & 它是运营资产权益，不是现金替代物或固定利息工具。 & 对每个项目单独检查资产权利、经营指标、维护资本开支、融资、治理、估值、扩募/解禁和成交能力。 \\
\bottomrule
\end{tabularx}
\sourcenote{S08--S14。规则为框架级摘要；具体交易前须核验当日有效规则、产品文件与账户权限。}
\end{exhibitbox}

\section{为什么结算与回转交易边界改变“止损”思维}

上交所现行交易规则列出当日回转交易的特定例外；普通买入股票的一般处置应按交收前不得卖出的规则设计，并须尊重具体品种、现行登记结算和交易规则 [S09]。深交所的现行规则亦应在每次实施时对照确认 [S12]。因此，普通 A 股策略不应把“跌到某个价格就立即卖掉”当作完整的风险控制。价格阈值的正确作用是触发一次\textbf{重新承保}：事实是否改变？熊市假设是否被击穿？若需要减仓，当前价格限制、可用流动性和账户规则下的可行路径是什么？

这并非反对纪律，而是把纪律从机械订单改成更难但更真实的三件事：事前仓位足够小、事前流动性足够好、事前失效条件足够明确。若一项投资只有在完美盘口或同日卖出条件下才安全，它在进入时就不安全。

\section{交易型开放式指数基金（ETF）、融资融券与“伪机构化”}

不同 ETF 的二级市场、申购赎回与回转安排取决于产品类别、产品文件和当期规则；即使某类产品存在申赎机制，也不等同于个人在二级市场拥有无摩擦的即时退出，且不能将某类 ETF 的机制泛化至权益、跨境、商品或 REIT [S13]。类似地，融资融券受标的范围、期限、保证金、担保物、维持担保和券商合同等条件限制 [S14]。本框架把它们都视为需要单独核验的产品/账户属性，而非一个可拿来放大私募收益的按钮。

个人投资者最常见的伪机构化行为是：用融资增加看似“高确定性”的仓位、用高分派产品替代现金、或以为更多代码就意味着更分散。这三者都可能让表面上的复杂度掩盖同一个真实风险因子：资金紧张、利率敏感、地产/租金、特定产业需求、或流动性抽离。

\section{公募不动产投资信托基金（REIT）：可研究的实物资产代理，但不是私募基础设施}

Brookfield 的所有者兼经营者模式启发的是研究顺序，而不是可照搬的持有结构 [S04]。上交所的公开募集不动产 REIT 业务办法及其发售上市指南提供相关制度语境 [S10--S11]；该办法衔接基础设施 REIT 指引与商业不动产 REIT 试点语境，实际适用范围仍须由每个产品的底层资产类别、现行规则和产品文件确认。每一个项目和产品仍需单独回答：

\begin{itemize}
  \item 现金流来自出租、车流、能源业务量、价格机制还是其他权利？其可持续性与合同/监管基础是什么？
  \item 经营 KPI 是出租率、加权平均租约期限、回款、车流、利用率、价格政策还是客户集中度？
  \item 维护资本开支、债务到期、再融资、扩募、解禁、关联交易和运营商激励如何影响可分配现金？
  \item 价格相对项目价值和未来可分配现金的含义是什么？市场成交能否支持预期的减仓路径？
\end{itemize}

只有这些问题有现时、项目级的披露答案时，REIT 才是“部分可迁移”的实物资产研究对象。历史高分派率本身不是资产质量、流动性或安全边际的证明。

\begin{riskbox}[实施前的强制暂停]
没有账户权限、当前产品合同、税费、买卖价差/深度、正常日成交能力、投资者风险承受能力和经批准的组合风险预算时，任何品种都不得从本报告的 \texttt{q=0} 状态升级。本备忘录提供研究逻辑，不授予交易资格。
\end{riskbox}
